\section{Sewing and local electrical observations}
\label{sec:inputs}

Electrical energy depends on every edge, including parallel edges.
Consequently, a description that approximates graph distances does not
by itself suffice for our purpose.  This section explains how the local
sewing construction of \cite{MetricMaps} retains exact graph pieces and
gives conditional laws for numerical observations on those pieces.
The observations may be extremal lengths or constrained energy minima;
we do not yet assert that they are bounded or identify their limits.
The construction has three parts.  We recover complete finite graph
pieces from contour cells, pass their numerical observations to
conditional limiting laws, and then read those laws from the fields
in the corresponding spatial regions.

Throughout this section,
\[
 \gamma=\sqrt2,\qquad Q=\frac2\gamma+\frac\gamma2=\frac3{\sqrt2},
 \qquad \chi=\frac1{\sqrt2}.
\]
All quantum surfaces use the field and boundary-length normalizations of
\cite{MetricMaps}.  A compact containment $K\Subset U$ means that $K$ is
compact and contained in the open coordinate domain $U$.  A buffer for
$K$ is an open neighborhood whose closure is compactly contained in $U$.

\subsection{The discrete graph and its continuum background}

Let $Z^{\mathrm d}=(L^{\mathrm d},R^{\mathrm d})$ be a bilateral walk
whose increments are independent and uniform on the four axial unit
vectors.  For integer times $i\le j$, set
\begin{equation}\label{inp:vertex-equivalence}
 i\sim_L j
 \quad\Longleftrightarrow\quad
 L^{\mathrm d}_i=L^{\mathrm d}_j
       =\min_{i\le k\le j}L^{\mathrm d}_k.
\end{equation}
The equivalence classes are the primal vertices.  Horizontal steps give
tree edges between consecutive classes, with the two traversals of a
tree edge counted as one edge.  A vertical up-step ending at $u$ and
the vertical down-step ending at $v>u$ are matched when
\begin{equation}\label{inp:step-match}
 R^{\mathrm d}_{u-1}=R^{\mathrm d}_v
      <\min_{u\le k<v}R^{\mathrm d}_k.
\end{equation}
Each matched pair contributes one further edge between $[u-1]_L$ and
$[v-1]_L$.  Distinct matched pairs remain distinct edges when their
endpoints agree.  In particular loops and parallel edges are retained.
Interchanging $L^{\mathrm d}$ and $R^{\mathrm d}$ gives the dual graph
and its complementary spanning tree.

For a finite interval of indices, the \emph{internal graph} uses the
same vertex relation, all horizontal steps in that interval, and the
vertical pairs whose two occurrences lie in the interval.  Its vertices
embed injectively into the full graph: the minimum test between two
retained times is the same in both constructions.  It need not contain
every full-graph edge between those vertices.

\begin{proposition}[Sewing and continuum background]\label{inp:sewing}
The following couplings and normalizations will be used.
\begin{enumerate}
\item A uniform rooted map--spanning-tree pair with $n$ edges is encoded
by the axial walk conditioned to remain in $\mathbb Z_{\ge0}^2$ and
return to zero after $2n$ steps.  The preceding rules recover every
primal and dual edge with its incidences.
\item Give the bilateral walk an independent Poisson clock of rate
$m$ and multiply its heights by $m^{-1/2}$.  As $m\to\infty$, its
interpolated contour converges locally uniformly in distribution to
$Z=(L,R)$, whose independent Brownian coordinates have variance
$1/2$ per unit time.  The limiting contour has the peanosphere coupling
with a $\sqrt2$-quantum cone $h$ and an independent whole-plane
space-filling $\mathrm{SLE}_8$, subsequently parameterized by $\mu_h$.
Write this curve as $\eta$ and normalize $\eta(0)=0$.
\item The imaginary field $\widehat h$ is independent of $h$ and is a
whole-plane GFF modulo $2\pi\chi\mathbb Z$, with uniform phase in
one period.  It determines the oriented curve before area
parameterization.  The fields, curve, and contour are retained with
their joint peanosphere law.
\item A finite interval of contour increments determines its abstract
decorated quantum surface, its restricted curve, and the quantum lengths
of its incoming and outgoing frontiers, modulo conformal coordinates.
The analogous statement holds for every symmetric interior interval
of the Brownian-excursion encoding of the ordinary unit-area sphere.
\end{enumerate}
\end{proposition}

The discrete assertion is Mullin's encoding \cite{Mullin}, in the form
of Theorem ``Mullin encoding and its infinite version'' in
\cite{MetricMaps}.  The continuum assertions are recorded in its
Theorem ``Peanosphere and interval surfaces'' and its imaginary-field
normalization.  Their principal continuum inputs are the mating of
trees theorems of \cite[Theorems~1.9 and~1.11]{DMS} and the
space-filling construction of \cite{IG4}.
The use of a Poisson clock makes the two coordinate jump processes
independent; index time divided by $m$ has the same limiting area time.

We recall the geometric information accompanying this coupling.  Two
continuum times have the same image under $\eta$ precisely when they
are equivalent under the relation generated by
\begin{equation}\label{inp:brownian-fibers}
 L_s=L_t=\inf_{[s,t]}L
 \quad\hbox{or}\quad
 R_s=R_t=\inf_{[s,t]}R,\qquad s\le t.
\end{equation}
Except for the common endpoint in the sphere encoding, every class has
at most three elements and belongs to one coordinate when it is
nonsingleton.  In a class $t_1<t_2<t_3$, the middle time is a strict
local minimum; distinct strict local minima have different heights.
At parameter $8$, the interval cells are closed topological disks.
The cone curve is proper, so inverse images of compact sets are bounded.
These facts, including their simultaneous versions on a countable
collection of intervals and independent Poisson cuts, are the
topological input accompanying Theorem ``Peanosphere and interval surfaces'' of
\cite{MetricMaps}; see also the identification description in
\cite[Section~1.3.1 and Lemma~8.15]{DMS}.

All subsequential arguments below can be placed in a coupling with
locally uniform contour convergence.  We attach the full continuum
background conditionally on the limiting contour, independently of
the additional \emph{raw} discrete observations given that contour.
Here raw means that the observation uses contour labels, without
projecting them through $\eta$.  This convention is compatible with
further extraction.  Indeed, for a bounded background test $F$, its
conditional expectation $k_F(Z)$ can be approximated in $L^1$ of the
fixed Brownian marginal by bounded continuous functions of $Z$.
Multiplying by any bounded raw test leaves the approximation error
uniformly bounded by the same $L^1$ error.  Joint weak convergence
therefore preserves the conditional sampling identity.  Spatially
projected observations will depend on both the raw data and the
background; we only make this conditional sampling assertion before
projection.

\subsection{Projections and planar topology}

Regard the map together with its dual as a cell decomposition of a
surface, subdivided into triangles at its corners.  Assign a primal or
dual vertex the point $\eta(t)$ corresponding to any of its visits.
All choices at one vertex, and all corners of an incident triangle,
have uniformly vanishing discrepancy.  Otherwise a subsequence of
visit times would limit to a discrete equality or matching relation
whose endpoints have distinct continuum images, contradicting
\eqref{inp:brownian-fibers}.  Join nearby assigned points by short
coordinate arcs and fill the triangles in small disks.  This gives
continuous projections $\pi_m$ with vanishing projected triangle mesh.
On the plane this construction is made in a bounded window containing
the region under consideration, after preserving its necessary stars
and capping the window.  This protection is obtained directly from
lower contour visits, as in the proof below; it does not use a local
instruction graph or any numerical-array limit.

\begin{proposition}[Topological comparison]\label{inp:topology}
In the contour coupling, the projections $\pi_m$ may be chosen so that
the following conclusions hold on the sphere and on every buffered
compact part of the plane.  Each conclusion is eventual for each
fixed choice of margins.
\begin{enumerate}
\item If the projections of two primal vertices approach each other,
they can be joined by primal paths with arbitrarily small projected
diameter.  This is uniform over all lifts in a fixed compact region.
The same statement holds on the dual graph and for surface paths.
In particular fixed tubes with clearance contain paths following them
on either graph.
\item All lifts of a sufficiently small closed disk or other compact
well can be included in a connected set over any fixed open
enlargement.  A loop projecting into a small disk can be filled over
a slightly larger disk, provided that both have the prescribed
clearance.
\item The projection has degree $+1$ or $-1$ in a compact spherical
completion.  Robustly transverse primal and dual tube crossings
intersect at an edge and its dual.  Projected circuits with nonzero
winding difference about two separated wells separate all the
corresponding lifted fibers.
\item Integer winding differences about two wells agree upstairs and
downstairs up to the common degree sign, when loops and wells have
fixed separation.  These statements retain repeated edge and face
incidences.
\end{enumerate}
\end{proposition}

\begin{proof}
The construction and the first three assertions are the construction
preceding Lemma ``Uniform topological joining'', that lemma,
Lemma ``Exterior connection and small-loop filling'', Proposition
``Degree and robust planar topology'', and Lemma ``Separation of
entire fibers'' in \cite{MetricMaps}.
We explain their applicability to both graphs and to the plane.

The hypotheses are uniform contour convergence, vanishing projected
cell mesh, and approximation of the limiting fiber generators by
discrete links with small projected connecting paths.  For a Brownian
excursion above a level, use last and first crossings of nearby levels
on the two sides of an interior time.  Their discrete analogues
approximate the chord.  At a triple, apply this on the two excursions
on either side of its middle minimum.  The fiber description above
therefore supplies every generator required for topological joining.
Interchanging the coordinates gives the same argument for the dual.
Applying planar separation to the two embedded graph paths gives the
mixed crossing assertion: a primal path and a dual path can meet only
through a primal edge and its dual, with the usual corner subdivision
making the intersection literal on the surface.

For a bounded plane region, choose a window containing all its visits
and its fixed additional star layers.  Both contour coordinates make
strictly lower visits on both sides of a sufficiently large window.
Lemma ``Protected capping of a bounded window'' in
\cite{MetricMaps} then supplies a spherical completion preserving
all these incidences.  Its continuum completion is absolutely
continuous with respect to a Brownian excursion, which supplies the
sphere quotient needed for the preceding topological results.

For the integer assertion, connect all lifts of each well away from
the tested loop, using the first two assertions.  A two-chain bounded
by that loop has constant multiplicity along each such connected
fiber.  Projection multiplies the difference of its multiplicities
at the two fibers by $\deg\pi_m=\pm1$.  This is exactly the equality
of winding differences.  Equivalently, one can use relative homology
after a simplicial perturbation inside the fixed margins.
\end{proof}

\subsection{Finite instruction graphs and numerical tests}

For a finite multigraph $G$, write
\begin{equation}\label{inp:graph-energy}
 \en_G(u)=\sum_{ab\in E(G)}(u(a)-u(b))^2.
\end{equation}
The sum has one term for each unoriented edge; loops contribute zero.
An energy test is the infimum of \eqref{inp:graph-energy} subject to
finitely specified value constraints at vertices of a specified
subgraph.  We allow $+\infty$ for an infeasible test.  An extremal-length
test for a family $\mathcal P$ of graph paths is
\begin{equation}\label{inp:el-test}
 \EL_G(\mathcal P)=
 \sup_{w:E(G)\to[0,\infty)}
 \frac{\bigl(\inf_{P\in\mathcal P}\sum_{e\in P}w(e)\bigr)^2}
      {\sum_{e\in E(G)}w(e)^2},
\end{equation}
with the usual conventions for empty families and constant paths.
Only nontrivial families will be used in the crossing arguments.

The graph on which one computes these quantities must be specified
before taking a limit.  We use the finite instruction graphs of
Lemma ``Finite-band reconstruction'' in \cite{MetricMaps}, recalled next.
Their inputs consist of finite strings and relative heights.  This
description will also show why the conditional-law proof applies to
our numerical tests.

\begin{definition}[Finite packages]\label{inp:packages}
A \emph{master block} consists of two increment strings on a time
interval $[0,\ell_b]$, starting at zero, together with a finite list
of internal cuts.  A finite package contains finitely many master
blocks.  For each coordinate $a\in\{L,R\}$ choose a forest on these
blocks and record its height differences
$o^a_{bd}=s^a_d-s^a_b$.  The values $s^a_b$ are thereby determined
within each forest component up to one additive constant.

A test chooses finitely many \emph{active windows}, each between
recorded cuts of a master block.  Each active window initially has its
own copy of the contour vertices.  Its starting heights and increment
strings are obtained by restriction of its master block.  Enable the
ordinary sewing inside each active window.  The additional
cross-window links are specified by the band rules below.  A package
can retain several tests on the same master strings.
\end{definition}

For an active window $b$, write $z^a_b$ for its increment string,
$s^a_b$ for its starting height, and $m^a_b=\min z^a_b$.
Its incoming and outgoing record arcs are the times satisfying,
respectively,
\[
 z^a_b(t)=\min_{0\le u\le t}z^a_b(u),
 \qquad
 z^a_b(t)=\min_{t\le u\le\ell_b}z^a_b(u).
\]
A band from window $b$ to window $d$ permits outgoing-to-incoming
links with common absolute height $H$ satisfying
\begin{equation}\label{inp:bands}
 \alpha\le H-q\le\beta,
 \qquad H\le s^a_w+\min_J z^a_w\quad(w,J\text{ in the cap list}).
\end{equation}
Here $\alpha<\beta$ are fixed rational parameters, $q$ is the height
at a recorded anchor cut, and each $J$ is a recorded subwindow.
The forest supplies every relative height appearing in the rule.
Each selected minimum is evaluated from the discrete string itself.

In the primal graph, an $L$-band identifies its permitted vertices.
An $R$-band is a rule for pairs of oriented step occurrences: an
outgoing up-step and an incoming down-step must carry the same unit
height strip, and the \emph{whole strip} must satisfy the band and cap
conditions.  Such a pair contributes an edge; it does not impose an
$R$-vertex identification.  The dual uses the exchanged convention.
If a tree edge is represented in several retained pieces it is counted
once after the corresponding vertex identifications: its unordered
endpoint pair identifies it, since a tree has no parallel edges.
Each additional matched-step edge is counted once for that matched
pair, even when another such pair has the same endpoints.

The same active windows and band instructions continue to define an
old test when cuts are added or master blocks are merged.  In the
merged string, old starting heights are evaluated at the old cuts;
they are no longer independent offset parameters.  A graph allowed
to use the entire merged interval is a separate new test.  In
particular an old test cannot be recovered by merely confining paths
in that larger graph, since the larger graph may use newly enabled
edges witnessed in a formerly omitted gap.

For each instruction list we retain a countable family of numerical
outputs.  A vertex is selected when one of its visits lies in a finite
union of intervals of relative time.  Endpoints can be recorded cuts
or rational fractions between cuts.  Such selectors specify the
subgraph, its ports, and any confinement condition.  Rational value
intervals at selected visits specify energy constraints; repeated
visits impose simultaneous constraints on the same vertex.  The
output library includes both graphs, finite joint lists of tests,
and tests recomputed on union graphs.  Inapplicable tests receive a
fixed default value.  Every nonnegative output is compactified in
$[0,\infty]$, for example by $x\mapsto x/(1+x)$.
We can retain outputs divided by any prescribed positive deterministic
factor $b_m$, including factors tending to zero.

This gives compact countable products of numerical observations.  The
master strings, durations, cuts, and offsets remain separate input
coordinates.  Along any sequence $m\to\infty$, one may therefore
extract joint limits of all the raw tests and the locally convergent
contours.  The issue is now to relate these formal tests to complete
pieces of the actual map.

\subsection{Exact cores and the selector sandwich}

Add independent Poisson cuts in area time, using nested refinements.
A cell is the image of a consecutive cut interval under $\eta$.
Only cells compactly contained in the current testing domain will be
used.  Their local data consist of their oriented traversals, quantum
area parameterizations, and tagged incoming and outgoing frontier
lengths, with repeated frontier occurrences retained.
A \emph{run} is a finite union of consecutive cells, viewed with its
concatenated traversal and tagged incoming and outgoing frontiers.
For any finite selection of cells, its maximal runs are the maximal
consecutive groups in that selection.

Here is how these local data provide the relative heights used in a
package.  Fix a run $b$ and one coordinate $a\in\{L,R\}$.  For an
initial segment of the run, the incoming frontier length is the drop
from its starting contour height to the minimum on that segment;
the outgoing length is the rise from that minimum to its ending
height.  Their difference is therefore the contour increment.
Reading these two tagged lengths at every rational relative area
time, and using continuity, recovers the whole increment path
$z_b^a$ from the local traversal.  Put
$m_b^a=\min z_b^a$ and $e_b^a=z_b^a(\ell_b)$.

The incoming frontier of the whole run has relative-height range
$[m_b^a,0]$, and the outgoing frontier has range $[m_b^a,e_b^a]$.
For a tagged point $p$ on either frontier, define its
\emph{relative height} $\vartheta_b^a(p)$ by starting with value
$m_b^a$ at the minimum endpoint and increasing by quantum boundary
length along that frontier toward its tip.  Thus the incoming tip
has height $0$ and the outgoing tip has height $e_b^a$.
This orientation is by increasing height, irrespective of the
orientation in which the contour traverses the frontier.  Equivalently,
if $p$ is represented by a record time $t$, then
$\vartheta_b^a(p)=z_b^a(t)$.

A \emph{directed shared seam} $b\to d$ is a common frontier subarc
of positive quantum length, tagged as outgoing for $b$ and incoming
for $d$, in the same coordinate.  At a paired point $p$ of this
seam, the two absolute heights agree.  Hence
\begin{equation}\label{inp:seam-offset}
 s_b^a+\vartheta_b^a(p)=s_d^a+\vartheta_d^a(p),
 \qquad
 s_d^a-s_b^a=\vartheta_b^a(p)-\vartheta_d^a(p).
\end{equation}
The right side is constant along the seam and is determined by the
two local tagged frontier parameterizations.  One may read it at an
interior point, so the construction does not require selecting or
assigning a height to a branch endpoint.  Taking a spanning forest
of such seams supplies exactly the offsets in
Definition~\ref{inp:packages}.  Their outgoing and incoming tags
also specify which linked run precedes the other.  A chronological
order between runs without a declared seam is unnecessary.

There is one endpoint configuration requiring an exact comparison.
Suppose three visits of one coordinate meet at height $H$, in three
disjoint runs $b,w,d$ in that order, with the middle visit a strict
local minimum.
The outer runs $b,d$ share a positive interval of seam heights below
$H$.  Above $H$, the two positive seam intervals are $b\to w$ and
$w\to d$; see Figure~\ref{inp:triple-figure}.  Formula
\eqref{inp:seam-offset} on these intervals recovers $s_w^a$ relative
to the outer starts.  The tag directions identify $w$ as the middle
run.  It lies between the two outer visits and contains the middle
contact.  The contour stays at least $H$ throughout that interval
by \eqref{inp:brownian-fibers}, and attains $H$ in $w$.  Hence
\begin{equation}\label{inp:exact-cap}
                       H=s_w^a+m_w^a.
\end{equation}
A band for the outer pair may extend above the continuum value
$H$, but it must include the cap in \eqref{inp:bands} evaluated
as $s_{w,m}^a+\min z_{w,m}^a$ at every discrete resolution.
It is this exact minimum, rather than a rounded approximation of
$H$, that excludes the outer-to-outer links above the branch.

\begin{figure}[htbp]
\centering
\begin{tikzpicture}[x=0.85cm,y=0.9cm,>=Stealth,
                    every node/.style={font=\small}]
 \fill[blue!7] (-0.15,2.1) rectangle (6.15,2.55);
 \fill[blue!7] (-0.15,0.4) rectangle (6.15,0.85);
 \draw[->,gray] (-1.7,0) -- (-1.7,3.1) node[above] {height};
 \draw[thick] (0,0.15) -- (0,2.95) node[above] {$b$};
 \draw[thick] (3,1.5) -- (3,2.95) node[above] {$w$};
 \draw[thick] (6,0.15) -- (6,2.95) node[above] {$d$};
 \draw[->,blue!65!black,thick] (0.08,2.3) -- (2.92,2.3);
 \draw[->,blue!65!black,thick] (3.08,2.3) -- (5.92,2.3);
 \draw[->,blue!65!black,thick] (0.08,0.62) -- (5.92,0.62);
 \node[above] at (1.5,2.3) {$b\to w$};
 \node[above] at (4.5,2.3) {$w\to d$};
 \node[below] at (3,0.62) {$b\to d$};
 \draw[densely dashed] (-1.9,1.5) -- (6.4,1.5);
 \node[left] at (-1.9,1.5) {$H$};
 \fill (3,1.5) circle (2pt);
 \node[anchor=west,fill=white,inner sep=2pt] at (3.25,1.73)
            {$s_w^a+m_w^a$};
 \draw[red!65!black,very thick] (-0.65,0.45) -- (-0.65,2.5);
 \draw[red!65!black] (-0.85,0.45) -- (-0.45,0.45);
 \draw[red!65!black] (-0.85,2.5) -- (-0.45,2.5);
 \node[anchor=east,rotate=90,red!65!black] at (-0.93,2.45)
            {uncapped band};
\end{tikzpicture}
\caption{A triple contact in height coordinates.  The shaded strips
indicate positive intervals of shared quantum-boundary heights;
they are not drawings of Euclidean seams.  The outer pair is linked
below $H$, while above $H$ each outer run is linked to the middle
run.  The middle run's exact discrete minimum caps the outer-pair
band.}
\label{inp:triple-figure}
\end{figure}

We use bands lying inside these shared height intervals with positive
margins at their effective endpoints, with two exceptions prescribed
by the record data.  At a natural lower record minimum the rational
band starts below the minimum and is clipped by the exact record
condition.  At an upper endpoint with the three-seam configuration,
the band is clipped by the exact selected minimum
\eqref{inp:exact-cap}.  These are the admissible band choices in
the proof of Lemma ``Finite-band reconstruction'' of
\cite{MetricMaps}.  All their tests use the local tagged frontiers,
relative increments, and quantum lengths just described.

\begin{samepage}
\begin{lemma}[Complete graph pieces in a buffer]\label{inp:complete-core}
Fix $K\Subset U$ in a coordinate chart.  Almost surely, there is a
finite prescription, selectable from the local cells in $U$ and their
frontier data, with the following property.  In a neighborhood of $K$,
its instruction graph agrees with the actual discrete graph, including
vertex equalities and all edge incidences, for all sufficiently large
resolutions in the coupling.  The same prescription can retain both
the primal and the dual graph.  One can require any fixed smaller
collar of this neighborhood before restricting the tests to it.

The prescriptions can be taken from the countable instruction scheme
above and chosen measurably using the fields and spatial cut marks
in $U$, with an arbitrarily small open buffer.  The assertion also
holds in the interior of a fixed interval cell, using only its
internal increment string.
\end{lemma}
\end{samepage}

\begin{proof}
We use the exact-incidence part of the proof of Lemma ``Finite-band
reconstruction'' in \cite{MetricMaps}.  We give the argument because
exactness, rather than an approximation of path costs, is essential
for energy tests.

Choose $r>0$ so that the closed $3r$-neighborhood of $K$ lies in $U$.
Properness bounds all its visit times in a finite window.  Refine the
Poisson cuts until every cell meeting this neighborhood has diameter
less than $r$.  Select all cells meeting the closed $2r$-neighborhood
of $K$ and merge consecutive selected cells into maximal runs.  These
runs are disjoint intervals; all visits sufficiently near $K$ lie in
them.  Their endpoints stay away from $K$, since otherwise the
neighboring cell would also have been selected.  Local traversal
data identify the selection: first inspect a slightly larger buffer
and require the portions around every visit to the smaller compact
to be flanked by marks while they remain in that buffer.  Properness,
continuity, and refinement make this condition hold eventually.
No global index for the retained runs is part of this inspection.

Consider one coordinate and two selected runs $b<d$ in chronological
order.  An outgoing point of $b$ and an incoming point of $d$ at
height $H$ are linked exactly when
\begin{equation}\label{inp:gap-test}
 H\le\min_{\text{between the two runs}} Z^a.
\end{equation}
The same assertion holds discretely, with the full-strip convention
for a matched step.  Split the intervening interval into selected
runs and omitted gaps.  If a limiting link projects near $K$, any
intervening time at its height represents that same continuum point.
It must therefore belong to a selected cell.  Thus every omitted-gap
minimum is strictly above the height of that link.  The relevant
endpoint pairs form a compact set, and there are only finitely many
gaps.  Their positive slack is consequently uniform: a sequence of
slacks tending to zero would limit to a contact in an omitted gap,
contrary to the selection.

An obstruction without slack can only be a minimum of a selected run.
At a triple contact, the three directed seams in
Figure~\ref{inp:triple-figure} identify the middle run and recover
its offset through \eqref{inp:seam-offset}.  Formula
\eqref{inp:exact-cap}, evaluated on the discrete string, supplies
the cap in \eqref{inp:bands}.  At a natural
lower seam endpoint, choose the lower rational band endpoint below
the record minimum; the exact incoming or outgoing record condition
then clips the band at that minimum.  This retains links at the
minimum as well as links at nearby lattice heights.

For every pair of runs, cover the compact set of relevant link heights
by finitely many such bands.  Choose each band to be admissible over
its entire effective height interval: that interval lies in a genuine
shared seam, its omitted-gap comparisons have positive slack, and
every allowed non-strict cap is the exact selected minimum just
described.  The natural record conditions provide the exact clipping
at lower seam endpoints.  Include the positive seam intervals
needed to read the relative heights, and retain a spanning forest
of these comparisons in each coordinate.  All needed offsets are
sums along that forest.  Seam directions specify the order of
linked runs and distinguish each middle cap; the instructions
do not need an order between unrelated runs.  Each band retains every discrete link
satisfying its rule.  Uniform contour convergence, the positive
omitted-gap slack, and the exact selected minima show two things:
every actual link near $K$ belongs to one of these bands eventually,
and every permitted link in those effective bands is an actual link
eventually.  In particular the latter assertion applies to any
permitted link with one endpoint over a smaller core collar: its
two limiting endpoints lie on the same shared-seam point and hence
remain inside the captured collar.  The
assertions hold simultaneously for all links, regardless of their
number.

Vertex equality is the same first-coordinate minimum test, so it is
captured as well.  Start a chain of enabled identifications at a
core representative.  Its first link is an actual identification
by the preceding band validity.  All visits to that actual vertex
have uniformly close projections, independently of the instruction
graph, and lie in the captured collar.  Inductively the next enabled
link is therefore actual as well and remains at that vertex.
Thus a peripheral chain cannot introduce an additional equality
at the core.  Horizontal tree traversals with the same identified
endpoint pair give one tree edge; no two distinct tree edges have
that pair.  Added edges are instead identified by their matched
step-occurrence pair, preserving distinct parallel edges.  Hence the
result is equality of graphs with incidences, not merely equality
of their path relations.  Applying the argument to the other
coordinate gives the dual conclusion.

All seam, cap, and positive-margin conditions are measurable in the
oriented local cells and frontier lengths.  Rational bands, integer
lower bounds on positive margins, and finite forest choices give a
countable selection rule.  Taking a further collar before the final
restriction gives the stated neighborhood of $K$.  Inside a fixed
interval cell, every visit to a compact subset of its interior lies
in that interval with a time margin; the same argument then uses
only its internal string.
\end{proof}

An arbitrary union of overlapping formal window copies need not be
an injective copy of a map subgraph.  The lemma asserts the existence
of \emph{complete} prescriptions, obtained from disjoint maximal
runs or from a freshly computed merged graph.  All comparisons to
the actual map below use these complete prescriptions.

\begin{lemma}[Spatial and value selectors]\label{inp:sandwich}
In the setting of Lemma~\ref{inp:complete-core}, suppose
$K\subset O$ with $K$ compact, $O$ open, and
$\overline O\Subset U$.  There is an interval selector on a complete
prescription which, eventually, includes every vertex projected into
$K$ and includes only vertices projected into $O$.  Several such
choices may be made together, for confinement and for each port set.

Let $u$ be continuous on a neighborhood of $\overline O$.  For any
two uniform value tolerances $0<\delta_-<\delta_+$, finite rational
interval constraints on selected visit labels can be placed between
the constraints $|v-u\circ\pi_m|\le\delta_-$ on a slightly larger
compact set and $|v-u\circ\pi_m|\le\delta_+$ on a smaller one.
All these comparisons hold simultaneously for every continuous
target through the same retained test library.
\end{lemma}

\begin{proof}
On the finitely many selected windows, the limiting parameterizations
are continuous.  Cover the inverse image of $K$ by finitely many
relative-time intervals with closures projecting inside $O$.
Their endpoints may be chosen from the prescribed dense countable
relative-time library.  Uniform projection accuracy and the complete
collar of Lemma~\ref{inp:complete-core} give the desired inclusions.
Select a vertex when any of its visits is in this union.  The same
argument applies to each port, while vanishing edge mesh allows an
induced-edge rule within the complete prescription.

For value constraints, refine the intervals until the oscillation of
$u\circ\eta$ on each is less than a fixed fraction of
$\delta_+-\delta_-$.  Choose a rational value interval between the
two allowed tolerances.  Impose it at every selected visit in that
interval, retaining all simultaneous constraints when a vertex has
several visits.  The strict tolerance gap and uniform projection
error give the two implications.  The construction chooses among
an already countable collection of tests; it does not add a new
exceptional null set for each target $u$.  Compact exhaustion and
rational positive tolerances therefore yield common measurable
versions for all continuous targets.
\end{proof}

The selector lemma compares variational values with strict margins.
It does not assert convergence for a prescribed sharp boundary.
For a crossing event, for example, choose smaller port and confinement
constraints and then relaxed constraints, both strictly inside the
desired geometry.  An intermediate selector produces actual paths
in the relaxed test, and any path family having the smaller clearance
can be included by a sufficiently relaxed selector.  Thus existence
of a strictly fitting test with finite limiting extremal length is
independent of the particular complete prescription.  Exhausted
energy values have the same property.  These conventions also make
readouts with a varying spatial center jointly measurable: all
membership and value comparisons use finite selectors with a strict
gap, followed by countable exhaustion.

\subsection{The conditional law of enlarged outputs}

We now identify the conditional law of the finite numerical arrays.
The proof is the finite-experiment argument of \cite{MetricMaps},
with the output vector enlarged before taking a subsequence.
The enlargement is harmless for a specific reason: every change
of density in that argument is a function of the inputs alone.

\begin{lemma}[Conditional transfer for numerical outputs]
\label{inp:kernel-transfer}
Consider finitely many packages as in Definition~\ref{inp:packages}.
Different packages use disjoint master blocks and disjoint offset
forests.  Sample their Poissonized strings independently; give their
durations, cuts, and offset vectors smooth positive densities, with
offsets discretized on the height lattice.  Retain any finite or
countable vector of compactified numerical tests of their instruction
graphs, including tests on both graphs.

Along a common subsequence the reference laws have disintegrations
\[
 \nu_r(\dd x_r)K_r(x_r,\dd y_r),
\]
where $x_r$ denotes the durations, cuts, increment strings, and
forest offsets, and $y_r$ denotes the numerical outputs.
Insert these blocks in the actual bilateral walk with positive
omitted gaps.  Conditional on the entire limiting contour and its
cuts, their output law is
\begin{equation}\label{inp:conditional-product}
             \bigotimes_r K_r(x_r,\dd y_r).
\end{equation}
The kernels are invariant under changes of positive reference
densities and under relabeling blocks or changing the forest basis.
They are compatible with jointly retained refinements and merged
tests.  These identities hold after selection from the countable
instruction library using the limiting contour and cuts.
\end{lemma}

\begin{proof}
This is the argument of Lemma ``Finite reference experiment and conditional
transfer'' in \cite{MetricMaps}.  We include its measure calculation
to specify its application to potentially discontinuous variational
outputs.

Write $\epsilon_m=m^{-1/2}$.  For one coordinate, order all selected
blocks chronologically as $1,\ldots,k$.  Let $s_i$ be their starts,
$e_i$ their net increments, and $g_i>0$ the intervening durations.
Complete the union of the package forests to a tree $T$ on these
blocks.  Give its extra offsets independent positive densities.
If the density on an edge $e$ is $\rho_e$, use lattice mass
\[
 r_{m,e}(a)=\int_{a-\epsilon_m/2}^{a+\epsilon_m/2}
                      \rho_e(u)\,\dd u,
 \qquad a\in\epsilon_m\mathbb Z.
\]
Let $p_t(j)$ be the transition mass of one coordinate of the
unscaled rate-one axial Poisson walk, so its variance rate is
$\sigma^2=1/2$.  The exact density of the actual relative starts
with respect to these reference offsets is
\begin{equation}\label{inp:density}
 W_m^a=
 \frac{\displaystyle\prod_{i=1}^{k-1}
       p_{mg_i}\bigl(\epsilon_m^{-1}(s_{i+1}-s_i-e_i)\bigr)}
      {\displaystyle\prod_{e\in T}r_{m,e}(a_e)}.
\end{equation}
Indeed, tree differences and consecutive differences are related
by an integer linear map with integer inverse, obtained by summing
along paths in $T$.  Its determinant is $\pm1$, so there is no
lattice multiplicity.  Multiply \eqref{inp:density} for the two
coordinates and include the density ratio for the durations and
cuts to obtain $W_m$.

On compact input sets with all gaps bounded below by a positive
constant, the local central limit theorem gives
\[
 \epsilon_m^{-1}p_{mt}(j)
 \longrightarrow
 (2\pi\sigma^2t)^{-1/2}
       \exp\{-x^2/(2\sigma^2t)\},
 \qquad x=\epsilon_m j,\quad \sigma^2=\tfrac12,
\]
uniformly in bounded $x$ and in such $t$.  The numerator and
denominator of \eqref{inp:density} have the same number of lattice
factors.  They cancel, leaving uniform convergence $W_m\to W$ to
a continuous finite function on these compact sets.

Conditional on the selected inputs, the omitted strings are
independent bridges and are independent of the numerical outputs.
Their rescaled laws converge uniformly on these input sets to
Brownian bridge laws.  For completeness, their finite-dimensional
densities are ratios of transition masses, so the preceding local
limit theorem applies.  On the first half of a bridge the density
relative to the free walk is bounded by
\[
 \frac{\sup_j p_{m(t-u)}(j)}{p_{mt}(j_0)}\le C,
 \qquad u\le t/2,
\]
when $t$ and $\epsilon_mj_0$ range in the chosen compact sets.
The numerator is at most $C\epsilon_m$ and the denominator at
least $c\epsilon_m$.  Free-walk tightness and the reversed argument
on the other half prove uniform bridge tightness.  The same applies
to a finite list of gaps and to exterior contour portions on a
bounded window.

The full independent reference experiments converge jointly to the
product of their limits, because their output spaces are compact
and their input marginals converge.  Let $a$ denote the completion
offsets and $z$ the omitted bridges.  Integrating bounded continuous
bridge tests first, the preceding uniform convergence gives the
limiting actual law
\begin{equation}\label{inp:weighted-law}
 W(x,a)\,\nu(\dd x,\dd a)\,
       \mathcal B(x,a;\dd z)\prod_rK_r(x_r,\dd y_r).
\end{equation}
The weight and the bridge law depend only on contour inputs.
The complete contour and cuts recover $(x,a,z)$, so disintegration
of \eqref{inp:weighted-law} gives \eqref{inp:conditional-product}.

To remove the compact restrictions, first keep all finitely many
limiting gaps positive, then bound endpoints and offsets, and
finally restrict the block paths to compact sets.  Under the actual
law the excluded probability tends to zero by finite-window contour
tightness.  If $P_m$ is the actual law and $Q_m$ the reference law,
the exact identity
\[
             \E_{Q_m}[W_m\1_{A^c}]=P_m(A^c)
\]
controls the omitted density mass.  No global bound on $W_m$ is
required.  Exhaustion, followed by a monotone-class argument, gives
the assertion for bounded measurable tests and for the entire
bilateral contour.

A change of reference density adds another input-only density
ratio, continuous in the limit on compact sets, and thus leaves
$K_r$ unchanged.  Forest bases are related by unimodular integer
maps.  The pushed-forward lattice cells and the directly discretized
cells have the same volume and vanishing diameter.  Their mass
ratios converge to one on compact sets; tightness makes the two
lattice laws converge in total variation.  This comparison remains
valid after applying a discontinuous graph-output map.  Absolute
time labels and the order of unrelated blocks occur only in
$W_m$, and disappear from the conditional output law.

For merged tests one must also check the old output.  Begin with
the old independent blocks, their offsets, and their old numerical
outputs.  Insert the formerly omitted positive-duration gaps by
the exact density and bridge construction.  The resulting string
has the new master-block law.  Recompute the old graph from its
unchanged active windows and bands, and compute the new union graph
as a separate output on that same string.  Formula
\eqref{inp:weighted-law}, applied before conditioning on the new
master string, proves the old conditional marginal.  It never
evaluates a previously chosen kernel on a forced-offset diagonal.
Windows with a common endpoint were already within one maximal
run, so no positive-density argument is applied to a zero-length
gap.  Jointly retaining finite vectors proves consistency for all
refinements and unions, and then for the countable library.

Every step used only compactness of the output coordinates and
their measurability from their own package.  Thus it applies to
the energy and extremal-length tests, with any deterministic
normalization.  It does not require continuity of those outputs
or preservation of conditional independence under arbitrary weak
limits.  Finally, after conditioning the identities hold for all
countably retained templates on one event; contour- and
cut-measurable selection from them is therefore permitted.
\end{proof}

\subsection{Spatial locality and changes of quantum units}

The preceding kernels are indexed by contour blocks.  To use them
in disjoint spatial regions, their inputs must be readable from
the fields in those regions.  This is where oriented cells and
relative frontier heights enter.

\begin{proposition}[Local numerical-array kernels]\label{inp:local-kernels}
On each open coordinate domain, take the limits of the numerical
tests on complete buffered prescriptions, with relative labels
localized to their cells.  Conditional on $(h,\widehat h)$, the
resulting local arrays admit kernels depending only on the fields
in an arbitrarily small open enlargement of the testing region.
Arrays using disjoint such enlargements are conditionally
independent given the complete fields.

One may retain the local Poisson marks.  Their conditional laws are
the corresponding Poisson processes with intensity proportional
to $\mu_h$, and their restrictions to disjoint domains are
independent.  The kernels are consistent under restriction,
common refinements, and compact exhaustion.
\end{proposition}

\begin{proof}
Proposition ``Background locality'' in \cite{MetricMaps}
states that $\widehat h$ in a buffer determines the oriented curve
segments through the testing region.  The restriction of $h$
gives their area parameterizations and their directed frontier
lengths.  These data determine the locally retained marked cells.
For a segment from time $s$ to time $t$, its coordinate increment
is the outgoing frontier length minus the incoming frontier
length of that coordinate, with tagged multiplicities retained:
\begin{equation}\label{inp:frontier-increments}
 Z^q_t-Z^q_s
   =\nu_h(\partial_q^{\mathrm{out}}\eta([s,t]))
      -\nu_h(\partial_q^{\mathrm{in}}\eta([s,t])),
 \qquad q\in\{L,R\}.
\end{equation}
Formula \eqref{inp:frontier-increments} reads
the increments without an absolute contour height.  Positive
shared seams read the forest offsets, and the three-seam
configuration in Lemma~\ref{inp:complete-core} reads the offsets
needed at a minimum cap.  Durations and relative cuts are local
area data.  Thus all four types of input in
Lemma~\ref{inp:kernel-transfer} are local.

For finitely many prescriptions in disjoint spatial domains,
merge within each domain.  Different domains use disjoint cells;
two runs from different domains cannot share a cut, since its
image would belong to both compactly contained cells.  Their
intervening chronological gaps are therefore positive even if
the runs interleave in time.  Lemma~\ref{inp:kernel-transfer}
gives their product law conditional on the contour and cuts.
Its invariance under labels and forest bases lets us enumerate
the cells locally without retaining a global cut index.

Adjoin the continuum background by the conditional sampling
convention above.  The product formula remains valid upon
conditioning additionally on the fields.  Its factors are now
measurable functions of the local fields and spatial marks,
by the preceding input reconstruction.  Conditional expectation
therefore gives the same product formula given the complete
fields and spatial marks.  This is exactly the argument of
Proposition ``Spatial product kernels'' in \cite{MetricMaps}, applied to the
enlarged output vector of Lemma~\ref{inp:kernel-transfer}.

The time cuts are independent of the continuum background.
Area parameterization and the Poisson
mapping theorem give their spatial conditional intensity
$\rho\mu_h$ for a layer of time intensity $\rho$.
Write $\Pi_i$ for the marks in the $i$th domain and
$K_i^{\mathrm{marks}}(h,\widehat h,\Pi_i;\cdot)$ for the
conditional output kernel already obtained.  Given the fields,
the marks have independent Poisson laws, and hence the output
law after removing them is
\[
 \bigotimes_i\int K_i^{\mathrm{marks}}
       (h,\widehat h,\pi_i;\cdot)\,
             \operatorname{Pois}_{\rho\mu_h|_{U_i}}(\dd\pi_i).
\]
This proves the fields-only assertion.  It uses no independence
between the marks and the numerical arrays, whose tests may
depend on those marks.  Alternatively, retain each $\Pi_i$
together with its array as the local output.
Finite joint consistency and the countable prescription library
then give the kernel on all tests.  Lemmas
\ref{inp:complete-core} and~\ref{inp:sandwich} identify its
buffered exhaustion with the corresponding observations from
the actual graph.  No additional conditioning on exterior
numerical outputs is used.
\end{proof}

\begin{proposition}[Coordinate and unit transport]\label{inp:transport}
The kernels in Proposition~\ref{inp:local-kernels} can be retained
jointly with the following transformations.
\begin{enumerate}
\item For a fixed conformal coordinate map $\phi$, pull back the
fields by
\[
 g\longmapsto g\circ\phi+Q\log|\phi'|,
 \qquad
 \widehat g\longmapsto\widehat g\circ\phi-\chi\arg\phi'.
\]
The numerical-array law is transported by the corresponding
spatial relabeling of its tests.
\item Fix $v=2^k$, $k\in\mathbb Z$, and put
$c_v=\gamma^{-1}\log v$.  Expressing the surface in the field
$g-c_v$ replaces area time by time divided by $v$ and replaces
rate $m$ by rate $mv$.  Electrical energies and extremal lengths
are unchanged by this relabeling.  If a numerical output is
divided by $b_m$, retain this same factor $b_m$ in the $mv$
experiment of the joint extraction.
\item Write $K_{(1)}$ for the original local law, including the
mark kernel, and $K_{(v)}$ for the law of the corresponding
transformed rate experiment.  Then
\begin{equation}\label{inp:unit-kernel}
 K_{(v)}(g,\widehat g)
     =\operatorname{Relabel}_v K_{(1)}(g+c_v,\widehat g).
\end{equation}
All fixed dyadic indices, their transformed templates, and the
required Poisson layers may be retained in one joint extraction.
\end{enumerate}
\end{proposition}

\begin{proof}
This is the finite-experiment identity underlying Proposition ``Intrinsic
transport and fixed quantum units'' in \cite{MetricMaps}.
A conformal change leaves quantum areas, frontier lengths,
directed seam roles, and their relative heights unchanged.
Hence it leaves the finite instruction graph unchanged and
only changes the projection of its labels.

Adding $c_v$ multiplies areas by $v$ and boundary lengths by
$\sqrt v$.  Thus the block inputs transform as
\begin{equation}\label{inp:block-units}
 \ell\longmapsto v\ell,
 \qquad z(t)\longmapsto\sqrt v\,z(t/v),
 \qquad o\longmapsto\sqrt v\,o.
\end{equation}
Cuts transform with time; band endpoints, their anchors, and
selected minima transform with height.  A rate-$m$ Poisson
string of duration $v\ell$ is, after this time change, a
rate-$mv$ string of duration $\ell$.  The height lattices agree
because $\sqrt v\,(mv)^{-1/2}=m^{-1/2}$.  Matching strips and
minimum comparisons are exactly the same integer tests.
Therefore the numerical costs on the resulting graph are
identical.  Smooth sampling densities may change, but
Lemma~\ref{inp:kernel-transfer} leaves the conditional kernels
unchanged.

The cuts must be transformed too.  At field $g+c_v$, the layer
corresponding to a fixed layer at field $g$ has time intensity
divided by $v$.  Its spatial intensity is then unchanged.
Choose a countable family of layers closed under these changes,
with independent layers and nested refinements as required.
Their joint input and output identities give
\eqref{inp:unit-kernel}.  A fixed field translation is locally
equivalent to the Gaussian reference law, as specified below;
this makes the identity valid under that fixed reference
marginal as well.  No estimate uniform in unbounded $v$ is
asserted or needed.
\end{proof}

For exhaustion of the plane we retain the joint cone normalization,
not only the deterministic-coordinate kernel identity.  In the
circle-average embedding define
\begin{equation}\label{inp:joint-zoom}
\begin{split}
 r_v&=\sup\{r>0:h_r(0)+Q\log r=c_v\},\\
 h^{(v)}&=h(r_v\,\cdot)+Q\log r_v-c_v,
 \qquad \widehat h^{(v)}=\widehat h(r_v\,\cdot),\\
 \eta^{(v)}(t)&=r_v^{-1}\eta(vt),
 \qquad Z^{(v)}(t)=v^{-1/2}Z(vt).
\end{split}
\end{equation}
The complete marginal has the original cone peanosphere law, by
the joint zoom in \cite{MetricMaps} and
\cite[Proposition~4.13(i)]{DMS}.  The positive dilation preserves
the independent stationary-phase imaginary-field law.  These
transformations are applied to the contour, raw arrays, and full
background together, before disintegration.  Thus field-dependent
reembedding does not require substituting a random map into
separately chosen versions of deterministic-coordinate kernels.
Moreover $r_{2^k}\to\infty$ as $k\to\infty$: the level $c_{2^k}$
diverges, whereas $h_r(0)+Q\log r$ is bounded above on compact
sets of radii and tends to $-\infty$ as $r\downarrow0$.
Ordinary patches inside these normalized disks therefore exhaust
the plane away from the cone point.

\begin{lemma}[Ordinary Gaussian reference patches]\label{inp:gaussian}
On a buffered patch compactly inside the cone's normalization
disk and away from its marked point, the actual real and
imaginary fields are locally equivalent to Gaussian reference
restrictions.  The real-field pinning circle can be fixed outside
all testing buffers.  This comparison preserves the actual
additive constant and can be made after any fixed affine change
of coordinates.

Adjoining any of the local kernels above preserves the same
field-only Radon--Nikodym derivative.  Consequently one fixed
density truncation transfers high-probability events uniformly
over all kernels attached to that reference marginal.  Fixed
additive real-field shifts and the fixed coordinate transforms
used in Proposition~\ref{inp:transport} preserve the corresponding
local null sets.
\end{lemma}

\begin{proof}
This is Lemma ``Ordinary cone patches and exterior Gaussian pins''
in \cite{MetricMaps}.  Inside the unit disk the circle-average
cone field has the law of the circle-average-zero whole-plane
GFF minus $\gamma\log|z|$.  On the patch, this deterministic
term has a smooth extension supported away from the marked
point and the pinning circle.  The Cameron--Martin theorem
and Gaussian domain comparison give local equivalence, with
the actual constant retained.  The same argument applies to
a fixed additive shift on the patch.  The imaginary field is
independent and is interpreted modulo its angular period;
its local restrictions have the corresponding Gaussian
equivalence and stationary-phase coordinate invariance.

If $Z=\dd\nu/\dd\mu$ is a field density and $K_i$ any of the
conditional kernels, then for every event $A$ and every $R>0$,
\begin{equation}\label{inp:density-truncation}
 (\nu\otimes K_i)(A)
 \le R(\mu\otimes K_i)(A)
      +\E_\mu[Z\1_{\{Z>R\}}].
\end{equation}
The last term does not depend on $i$.  This proves the uniform
probability comparison; local equivalence gives both directions
for probability-one events.  A test domain containing the full
pinning circle is excluded by the stated buffer condition.
\end{proof}

\begin{lemma}[Primal--dual reflection]\label{inp:reflection}
The joint local-array laws are invariant under exchanging the
primal and dual graphs and reflecting their coordinate tests.
They are also covariant under fixed rigid motions, with the
natural transformations of the two fields.
\end{lemma}

\begin{proof}
The discrete assertion follows by exchanging the two independent
coordinate strings and their forests.  In the continuum,
spatial reflection combined with negation of the reflected
imaginary field exchanges the two angle trees and the two
sides of the curve.  The space-filling ordering is preserved:
it can be read either from the $+\pi/2$ tree with right merges
or from the $-\pi/2$ tree with left merges, as in
\cite[Section~1.2.3, Figure~1.16]{IG4}.  Both descriptions are
exchanged by this operation, so no reversal of area time is
involved.  The cone field and the independent stationary-phase
imaginary field have the required reflected law.  Boundary
lengths exchange $L$ and $R$, and conditional sampling of the
background therefore gives the same operation on the joint
raw-array limit.  The intrinsic kernels and the selector
construction then transport the spatial tests.  Fixed rotations
and translations follow from coordinate covariance and local
field equivalence on buffered ordinary patches.  A countable
group of such transformations can be retained simultaneously.
\end{proof}

\begin{proposition}[Local-array interface]\label{inp:array-interface}
From every sequence of resolutions tending to infinity one can
extract, jointly with the continuum peanosphere background,
all compactified numerical tests described above, on both
graphs and at every fixed dyadic area offset.  Prescribed
deterministic normalizing factors may be retained even when
they tend to zero.  In this extraction:
\begin{enumerate}
\item complete buffered graph pieces and the selector sandwich
of Lemmas~\ref{inp:complete-core} and~\ref{inp:sandwich} identify
local tests with actual graph observations;
\item Proposition~\ref{inp:local-kernels} supplies buffer-local
conditional kernels and product laws in disjoint domains;
\item Proposition~\ref{inp:transport} and
Lemmas~\ref{inp:gaussian}--\ref{inp:reflection} supply their
coordinate, unit, Gaussian-reference, and primal--dual
comparisons;
\item these properties persist under further joint extraction
and under countable exhaustion with strict margins.
\end{enumerate}
\end{proposition}

\begin{proof}
All numerical coordinates belong to compact spaces, and the
input contours are tight on bounded windows.  Include the
countable transformed instruction lists and mark layers before
the diagonal extraction.  The preceding results apply to
every finite collection.  Their exact density and bridge
calculation passes to any further limit by compact truncation;
the continuum background passes by the $L^1$ conditional
sampling argument.  Countable intersection and the selector
exhaustions give the stated simultaneous version.
\end{proof}

The interface provides laws and exact local observability.
The bounds which prevent degeneration of the electrical tests,
and the identification of their variational limits, are proved
next.
